\documentclass[10pt,a4paper]{amsart}
\usepackage[margin=19mm]{geometry}
\usepackage{amsmath,amssymb,mathtools}
\usepackage[scr=rsfs,cal=euler]{mathalfa}
\usepackage{xfrac}
\usepackage[unicode,hidelinks,bookmarks=false]{hyperref}
\newcommand{\ie}{i.e.\ }
\newcommand{\cf}{cf.\ }
\newcommand{\resp}{respectively\ }
\newcommand{\Subsection}{\subsection}
\providecommand{\nobreakdash}{\nobreak\hbox{-}}
\setlength{\emergencystretch}{2em}
\multlinegap0pt
\usepackage{amssymb}
\usepackage{float}
\usepackage{mathtools}
\usepackage{cancel}
\newtheorem{lem}{Lemma}
\def \Ga{\Gamma}
\def \H{\mathbb{H}}
\def \K{\mathcal{K}}
\def \S{\mathbb{S}}
\def \T{\operatorname{T}}
\def \es{\emptyset}
\def \k{\Bbbk}
\newcommand{\lr}[1]{\langle #1 \rangle}
\def \sp{\operatorname{span}}
\title{Superalgebras and Algebras with Involution: Classifying Cubic Codimension Sequences}
\author{Yan-Hong Bao$^*$, Jiang-Nan Xu, Yuan-Feng Zhang}
\date{Source: arXiv:2608.15612v1 (CC BY 4.0)}
\begin{document}
\maketitle

\noindent\emph{Source and licence.} Superalgebras and Algebras with Involution: Classifying Cubic Codimension Sequences. Version arXiv:2608.15612v1 (CC BY 4.0), distributed by arXiv under the Creative Commons Attribution 4.0 International licence, http://creativecommons.org/licenses/by/4.0/, as recorded in the arXiv licence metadata for this exact version and shown on the abstract page https://arxiv.org/abs/2608.15612v1. Full licence notice: \texttt{licenses/arxiv-2608.15612-CC-BY-4.0.txt}.\par\smallskip

\noindent\textit{Editorial scope.} Counted unit: the article's lem lem2.12 (source0.tex lines 645--657). The article's complete original proof of it is delivered with it (source0.tex lines 659--743, one \texttt{proof} environment of 85 lines). No mathematics is written, repaired or re-derived: the statement and the proof are the article's own lines, in the article's own notation, and no retained line is substituted or re-flowed.  No further source-local context is needed: every cross-reference of every retained line resolves inside the two delivered spans.  Nothing was omitted from any retained span, so the package carries no omission record.  No retained line contains a citation, so the wrapper carries no bibliography and no bibliography entry.  No retained line contains a figure environment, an image-inclusion command or drawing code, so no image is delivered and the package declares no asset dependency.  Editorial formatting only, in addition: an AMS \texttt{amsart} preamble that restates the article's own declarations and macros used by the retained lines, namely the environments \texttt{lem} on separate counters, together with the standard packages those lines need.  The retained environments print in this document as Lemma 1 in order of appearance, whereas the article prints the same items with the numbering of its own full document.  The complete native source of the article as distributed in the arXiv e-print for this exact version is bundled unchanged as \texttt{sources/207765/source0.tex}.
\begin{lem}\label{lem2.12}
	Let $P_2$ be an irreducible submodule of $\Ga_2^*= V_{(1^2), \es} \oplus  V_{(1),(1)} \oplus V_{\es,(2)} \oplus V_{\es,(1^2)}.$
	\begin{enumerate}
		\item[(1)]
		If $P_2=V_{(1^2), \es}$, then $\chi({\lr{V_{(1^2), \es}}_{\T^*} \cap \Ga_3^*}) = \chi_{(2,1),\es} +\chi_{(2),(1)} + 2\chi_{(1^2),(1)} + \chi_{(1),(2)}$.
		\item[(2)]
		If $P_2=V_{(1), (1)}$, then $\chi({\lr{V_{(1), (1)}}_{\T^*} \cap \Ga_3^*}) = \chi_{(2,1),\es} + \chi_{(2),(1)} + 2\chi_{(1^2),(1)} + 2\chi_{(1),(2)} + 2\chi_{(1),(1^2)} + \chi_{\es, (2,1)}$.
		\item[(3)]
		If $P_2=V_{\es, (2)}$, then $\chi({\lr{V_{\es, (2)}}_{\T^*} \cap \Ga_3^*}) = \chi_{(1^2),(1)} + 2\chi_{(1),(2)} + \chi_{(1),(1^2)} + \chi_{\es, (3)} + \chi_{\es, (1^3)} + 2\chi_{\es, (2,1)}$.
		\item[(4)]
		If $P_2=V_{\es, (1^2)}$, then $\chi({\lr{V_{\es, (1^2)}}_{\T^*} \cap \Ga_3^*}) = \chi_{(1^2),(1)} + \chi_{(1),(2)} + 2\chi_{(1),(1^2)} + \chi_{\es, (1^3)} + 2\chi_{\es, (2,1)}$.
	\end{enumerate}
\end{lem}	

\begin{proof}
	Recall that bases for $\K^+ \cap V_2^*$ and $\K^- \cap V_2^*$ are given respectively by 
	\[
	\{y_1 \circ y_2, z_1 \circ z_2, [y_1,z_2], [z_1,y_2]\} \quad \text{and} \quad \{[y_1,y_2], [z_1,z_2], y_1 \circ z_2, z_1 \circ y_2\}.
	\]
	
	(1) Assume $P_2 = V_{(1^2), \es} = \sp_\k \{[y_2, y_1] \}$. Using the above bases of $\K^+ \cap V_2^*$ and $\K^- \cap V_2^*$, the left $\k\H_3$-module $\lr{V_{(1^2), \es}}_{\T^*} \cap V_3^*$ is generated by elements of the following forms:
	\[
	[y \circ y, y], \quad y[y,y], \quad [y,y]y, \qquad [[z,y],y], \quad z[y,y], \quad [y,y]z, \qquad [z \circ z, y].
	\]
	
	\textbf{Case 1.} Bipartition of type $(3,0)$. Combining the identity $[y_3,y_1,y_2]=[y_3,y_1]y_2-y_2[y_3,y_1]$ with the relevant Specht basis elements $[y_3,y_1,y_2]$ and $[y_3,y_2,y_1]$ of $\Ga_3^*$, we see that the $\H_3$-character $\chi\big(\lr{V_{(1^2), \es}}_{\T^*} \cap \Ga_3^*\big)$ contains exactly one copy of $\chi_{(2,1),\es}$.
	
	\textbf{Case 2.} Bipartition of type $(2,1)$. By the representation theory of symmetric groups it suffices to consider the $\S_2\times\S_1$-action. The relevant Specht basis elements of $\Ga_3^*$ are $[z_3,y_1,y_2]$, $[z_3,y_2,y_1]$ and $z_3[y_2,y_1]$, all of which are readily obtained from generators $[z,y,y]$ and $z[y,y]$. The transposition $(12)$ acts by
	\begin{align*}
		(12)\cdot([z_3,y_1,y_2]+[z_3,y_2,y_1])&=([z_3,y_1,y_2]+[z_3,y_2,y_1]),\\
		(12)\cdot([z_3,y_1,y_2]-[z_3,y_2,y_1])&=-([z_3,y_1,y_2]-[z_3,y_2,y_1]),\\
		(12)\cdot z_3[y_2,y_1]&=-z_3[y_2,y_1].
	\end{align*}
	Hence, the $\H_3$-character $\chi( \lr{V_{(1^2), \es}}_{\T^*} \cap \Ga_3^* )$ contains one copy of $\chi_{(2),(1)}$ and two copies of $\chi_{(1^2),(1)}$.	
	
	\textbf{Case 3.} Bipartition of type $(1,2)$. Here it is enough to consider the $\S_1\times\S_2$-action. The relevant Specht basis elements of $\Ga_3^*$ are $[z_3,y_1,z_2]$, $[z_3,z_2,y_1]$, $z_2[z_3,y_1]$ and $z_3[z_2,y_1]$. The generator $[z\circ z,y]$ produces only $[z_2\circ z_3,y_1]$. Note that
	\[
	[z_2 \circ z_3,y_1]=2[z_3,y_1,z_2]-[z_3,z_2,y_1]+2z_2[z_3,y_1]+2z_3[z_2,y_1] = f_{(1),(2)} \in \Ga_3^*,
	\]
	Since $(23)\cdot [z_2 \circ z_3,y_1] = [z_2 \circ z_3,y_1]$, the $\H_3$-character $\chi({\lr{V_{(1^2), \es}}_{\T^*} \cap \Ga_3^*})$ contains one copy of $\chi_{(1),(2)}$.
	
	Combining the three cases we obtain
	\[
	\chi({\lr{V_{(1^2), \es}}_{\T^*} \cap \Ga_3^*}) = \chi_{(2,1),\es} +\chi_{(2),(1)} + 2\chi_{(1^2),(1)} + \chi_{(1),(2)}.
	\] 
	
	In parts (2), (3) and (4), we only detail the bipartitions that appear with multiplicity $2$ in $\chi(\Ga_3^*)$ but with multiplicity $1$ in the corresponding $\chi(\langle P_2\rangle_{\T^*} \cap \Ga_3^*)$. The remaining bipartitions follow by the same arguments as in (1).
	
	(2) Assume $P_2 = V_{(1), (1)} = \sp_\k \{[z_2,y_1], [y_2,z_1] \}$. The $\k\H_3$-module $\lr{V_{(1),(1)}}_{\T^*} \cap V_3^*$ is generated by
	\[
	[[y,y],y],\enspace\enspace [z\circ y,y],\enspace [z,y\circ y],\enspace y[z,y],\enspace [z,y]y,\enspace\enspace [[z,z],y],\enspace [z,[z,y]],\enspace  z[z,y],\enspace  [z,y]z,\enspace\enspace [z,z\circ z].
	\] 
	We concentrate on the bipartition $(\es, (2,1))$, for which the only relevant generator is $[z,z\circ z]$. Observe that $[z_1,z_2\circ z_3]$ and $(12)\cdot [z_1,z_2\circ z_3]$ are linearly independent. Under the involution $*$,
	\begin{align*}
	*([z_1,z_2\circ z_3])=[z_1,z_2\circ z_3]=-\frac{2}{3}g_{\es,(2,1)}^1+\frac{1}{3}g_{\es,(2,1)}^2,\\
	*((12)\cdot [z_1,z_2\circ z_3])=(12)\cdot [z_1,z_2\circ z_3]=\frac{1}{3}g_{\es,(2,1)}^1-\frac{2}{3}g_{\es,(2,1)}^2.
	\end{align*}
	Consequently, the $\H_3$-character $\chi(\lr{V_{(1),(1)}}_{\T^*} \cap \Ga_3^*)$ contains one copy of $\chi_{\es,(2,1)}$.
	
	(3) Assume $P_2 = V_{\es, (2)} = \sp_\k \{z_1 \circ z_2 \}$. The $\k\H_3$-module $\lr{V_{\es, (2)}}_{\T^*} \cap V_3^*$ is generated by 
	\[
	[y,y]\circ z,\qquad (y\circ z)\circ z,\quad y(z\circ z),\quad (z\circ z)y,\qquad [z,z]\circ z,\quad z(z\circ z),\quad (z\circ z)z.
	\] 
	We first consider the bipartition $((1^2), (1))$. The only relevant element is $[y_1,y_2] \circ z_3$. Since
	\[
	[y_1,y_2]\circ z_3=-([z_3,y_1,y_2]-[z_3,y_2,y_1])-2z_3[y_2,y_1]= -g_{(1^2),(1)} \in \Ga_3^*,
	\]
	the $\H_3$-character $\chi(\lr{V_{\es, (2)}}_{\T^*} \cap \Ga_3^*)$ contains one copy of $\chi_{(1^2),(1)}$.
	
	We now turn to the bipartition $((1), (1^2))$. The three forms $(y\circ z)\circ z$, $y(z\circ z)$ and $(z\circ z)y$ yield the following maximal linearly independent set:
	\[
	(y_1\circ z_2)\circ z_3, \quad (y_1\circ z_3)\circ z_2, \quad y_1(z_2\circ z_3), \quad (z_2\circ z_3)y_1.
	\]
	Examining the action of the transposition (23) shows that
	\[
	(y_1\circ z_2)\circ z_3-(y_1\circ z_3)\circ z_2=[z_3,z_2,y_1] = g_{(1),(1^2)} \in \Ga_3^*
	\]
	Thus $\chi(\lr{V_{\es, (2)}}_{\T^*} \cap \Ga_3^*)$ also  contains one copy of $\chi_{(1),(1^2)}$.
	
	(4) Assume $P_2 = V_{\es, (1^2)} = \sp_\k \{ [z_1,z_2] \}$. The $\k\H_3$-module $\lr{V_{\es, (1^2)}}_{\T^*} \cap V_3^*$ is generated by
	\[
	[[y,y],z],\qquad [y\circ z,z],\quad y[z,z],\quad [z,z]y,\qquad [[z,z],z],\quad z[z,z],\quad [z,z]z
	\]
	We first consider the bipartition $((1^2), (1))$. The only relevant element is $[[y_1,y_2],z_3]$. Since
	\[
	[y_1,y_2,z_3]=[z_3,y_2,y_1]-[z_3,y_1,y_2] = -f_{(1^2),(1)} \in \Ga_3^*
	\]
	the $\H_3$-character $\chi(\lr{V_{\es, (1^2)}}_{\T^*} \cap \Ga_3^*)$ contains one copy of $\chi_{(1^2),(1)}$.
	
	We now turn to the bipartition $((1), (2))$. The three forms $[y\circ z,z]$, $y[z,z]$ and $[z,z]y$ yield the following maximal linearly independent set:
	\[
	[y_1\circ z_2,z_3], \quad [y_1\circ z_3,z_2], \quad y_1[z_3,z_2], \quad [z_3,z_2]y_1.
	\]
	Examining the action of the transposition (23) shows that
	\[
	[y_1\circ z_2,z_3]+[y_1\circ z_3,z_2]=-(2[z_3,y_1,z_2]-[z_3,z_2,y_1])-2(z_2[z_3,y_1]+z_3[z_2,y_1]) = -f_{(1),(2)} \in \Ga_3^*,
	\]
	Thus $\chi(\lr{V_{\es, (1^2)}}_{\T^*} \cap \Ga_3^*)$ also  contains one copy of $\chi_{(1),(2)}$.
\end{proof}	

\end{document}
